% preamble.tex — shared packages, colors, and styles.
% Loaded once from main.tex via \input. Do not add \documentclass here.

\usepackage[T1]{fontenc}
\usepackage{lmodern}
\usepackage[utf8]{inputenc}
\usepackage{amsmath,amsfonts,amssymb}
\usepackage{geometry}
\geometry{a4paper, margin=1in}
\usepackage{xcolor}
\usepackage{listings}
\usepackage{hyperref}
\usepackage[most]{tcolorbox}
\usepackage{titlesec}
\usepackage{graphicx}
\usepackage{parskip}
\usepackage{booktabs}
\usepackage{float}
\usepackage{makeidx}
\makeindex
\usepackage{tikz}
\usetikzlibrary{positioning, arrows.meta, shapes.geometric, calc, fit, backgrounds}
\usepackage{pgfplots}
\pgfplotsset{compat=1.18}

% Colors
\definecolor{codegreen}{rgb}{0,0.6,0}
\definecolor{codegray}{rgb}{0.5,0.5,0.5}
\definecolor{codepurple}{rgb}{0.58,0,0.82}
\definecolor{backcolour}{rgb}{0.96,0.96,0.96}
\definecolor{primary}{rgb}{0.0, 0.3, 0.6}

% Shared TikZ styles for chapter diagrams (reuse preamble colors)
\tikzset{
    stage/.style={
        rectangle, rounded corners, draw=primary, thick,
        fill=backcolour, text=black, align=center,
        minimum width=2.4cm, minimum height=1cm, font=\small
    },
    llmstage/.style={
        rectangle, rounded corners, draw=magenta, thick,
        fill=magenta!10, text=black, align=center,
        minimum width=2.4cm, minimum height=1cm, font=\small
    },
    store/.style={
        cylinder, shape aspect=0.25, draw=primary, thick,
        fill=backcolour, shape border rotate=90,
        align=center, minimum width=1.8cm, minimum height=1.3cm, font=\small
    },
    tool/.style={
        rectangle, draw=codegreen!60!black, thick, fill=green!6,
        align=center, minimum width=1.8cm, minimum height=0.9cm, font=\footnotesize
    },
    doc/.style={
        rectangle, draw=codegray, thick, fill=white,
        align=center, minimum width=1.6cm, minimum height=0.8cm, font=\footnotesize
    },
    note/.style={font=\footnotesize\itshape, text=codegray},
    arrow/.style={-{Stealth[length=2.5mm]}, draw=primary, thick},
    biarrow/.style={{Stealth[length=2mm]}-{Stealth[length=2mm]}, draw=codegray, thick}
}

\hypersetup{
    colorlinks=true,
    linkcolor=primary,
    citecolor=primary,
    filecolor=magenta,
    urlcolor=primary,
    bookmarksnumbered=true,
    bookmarksopen=true,
    pdftitle={Azure Master Data Engineer},
    pdfauthor={PhD Aned Esquerra Arguelles},
    pdfsubject={A Comprehensive, Hands-On Academic Guide to Data Engineering on Microsoft Azure},
    pdfkeywords={Azure Data Lake Storage Gen2, ADLS Gen2, Azure Synapse Analytics, Azure Data Factory, Azure Databricks, Cosmos DB, Event Hubs, Delta Lake, data lakehouse, data governance, data engineering},
}

% Code Listing Style
\lstdefinestyle{pystyle}{
    backgroundcolor=\color{backcolour},
    commentstyle=\color{codegreen},
    keywordstyle=\color{magenta},
    numberstyle=\tiny\color{codegray},
    stringstyle=\color{codepurple},
    basicstyle=\ttfamily\footnotesize,
    breakatwhitespace=false,
    breaklines=true,
    captionpos=b,
    keepspaces=true,
    numbers=left,
    numbersep=5pt,
    showspaces=false,
    showstringspaces=false,
    showtabs=false,
    tabsize=4,
    language=Python
}
\lstset{style=pystyle}

\lstdefinelanguage{json}{
    morestring=[b]",
    morecomment=[l]{//},
    literate=
     *{0}{{{\color{magenta}0}}}{1}
      {1}{{{\color{magenta}1}}}{1}
      {2}{{{\color{magenta}2}}}{1}
      {3}{{{\color{magenta}3}}}{1}
      {4}{{{\color{magenta}4}}}{1}
      {5}{{{\color{magenta}5}}}{1}
      {6}{{{\color{magenta}6}}}{1}
      {7}{{{\color{magenta}7}}}{1}
      {8}{{{\color{magenta}8}}}{1}
      {9}{{{\color{magenta}9}}}{1}
}
\lstdefinelanguage{AzureCLI}{
    language=bash,
    morekeywords={az,group,create,storage,account,container,policy,management-policy,update,show,login,extension,add,role,assignment,fs,directory,file,upload,download,exists},
    sensitive=true
}

% Chapter Styling
\titleformat{\chapter}[display]
  {\normalfont\bfseries\color{primary}}{\chaptertitlename\ \thechapter}{20pt}{\Huge}

% Smart cross-references (loaded after hyperref, per cleveref requirements).
\IfFileExists{cleveref.sty}{%
    \usepackage{cleveref}%
}{%
    % Portable fallback when cleveref is unavailable: keep \cref/\Cref working.
    \providecommand{\cref}[1]{\ref{#1}}%
    \providecommand{\Cref}[1]{\ref{#1}}%
    \providecommand{\crefrange}[2]{\ref{#1}--\ref{#2}}%
    \providecommand{\cpageref}[1]{\pageref{#1}}%
}

% ---------------------------------------------------------------------------
% Pedagogical callout boxes (reuse across every chapter for a consistent,
% scannable learning experience). All build on the shared color palette.
% ---------------------------------------------------------------------------
\newtcolorbox{objectives}{
    breakable, colback=primary!5, colframe=primary, arc=2pt,
    left=6pt, right=6pt, top=4pt, bottom=4pt,
    title={\textbf{Learning Objectives}}, fonttitle=\bfseries\color{white}
}
\newtcolorbox{keyterms}{
    breakable, colback=backcolour, colframe=codegray, arc=2pt,
    left=6pt, right=6pt, top=4pt, bottom=4pt,
    title={\textbf{Key Terms}}, fonttitle=\bfseries\color{white}
}
\newtcolorbox{definitionbox}{
    breakable, colback=backcolour, colframe=primary, arc=2pt,
    left=6pt, right=6pt, top=4pt, bottom=4pt,
    title={\textbf{Definition}}, fonttitle=\bfseries\color{white}
}
\newtcolorbox{notebox}{
    breakable, colback=blue!4, colframe=blue!45!black, arc=2pt,
    left=6pt, right=6pt, top=4pt, bottom=4pt,
    title={\textbf{Note}}, fonttitle=\bfseries\color{white}
}
\newtcolorbox{tipbox}{
    breakable, colback=green!4, colframe=green!45!black, arc=2pt,
    left=6pt, right=6pt, top=4pt, bottom=4pt,
    title={\textbf{Tip}}, fonttitle=\bfseries\color{white}
}
\newtcolorbox{pitfall}{
    breakable, colback=red!4, colframe=red!60!black, arc=2pt,
    left=6pt, right=6pt, top=4pt, bottom=4pt,
    title={\textbf{Common Pitfall}}, fonttitle=\bfseries\color{white}
}

% Cross-cloud equivalence callout (used to align concepts across the six platform books).
\newtcolorbox{crosscloud}{
    breakable, colback=teal!4, colframe=teal!55!black, arc=2pt,
    left=6pt, right=6pt, top=4pt, bottom=4pt,
    title={\textbf{Cross-Cloud Equivalence}}, fonttitle=\bfseries\color{white}
}

% Capstone project callout (opens the end-of-chapter hands-on capstone).
\newtcolorbox{capstonebox}{
    breakable, colback=red!3, colframe=red!55!black, arc=2pt,
    left=6pt, right=6pt, top=4pt, bottom=4pt,
    title={\textbf{Capstone Project}}, fonttitle=\bfseries\color{white}
}
